% Teil II: der durchgerechnete Weg.

\subsection*{Gegeben und gesucht}
(a) $\sum_{k=0}^{n}\binom{n}{k}(-2)^k$ auswerten; (b) DNA-Replikation mit
$n$ Basen, Mutationswahrscheinlichkeit $p$ je Base: keine bzw. genau $k$
Mutationen, und $\sum\binom{n}{k}p^k(1-p)^{n-k}=1$.

\subsection*{a) Wert der Summe}
Nach dem binomischen Lehrsatz mit $a=-2$, $b=1$:
\begin{equation}
  \sum_{k=0}^{n}\binom{n}{k}(-2)^k=\bigl((-2)+1\bigr)^n=(-1)^n .
\end{equation}
\begin{ergebnis}
$1$ für gerade $n$, $-1$ für ungerade $n$.
\end{ergebnis}

\subsection*{b) Mutationen}
i) Keine Mutation ($k=0$):
\begin{equation}
  \mathrm P=\binom{n}{0}p^0(1-p)^n=(1-p)^n .
\end{equation}

ii) Genau $k$ Mutationen:
\begin{equation}
  \mathrm P(\text{genau }k)=\binom{n}{k}p^k(1-p)^{n-k}.
\end{equation}

iii) Mit dem Lehrsatz für $a=p$, $b=1-p$:
\begin{equation}
  \sum_{k=0}^{n}\binom{n}{k}p^k(1-p)^{n-k}=(p+(1-p))^n=1 .
\end{equation}

\begin{ergebnis}
\begin{align*}
  \text{i)}\ &(1-p)^n,\\
  \text{ii)}\ &\binom{n}{k}p^k(1-p)^{n-k},\\
  \text{iii)}\ &\sum_{k=0}^{n}\binom{n}{k}p^k(1-p)^{n-k}=1,\ \text{also}\\
   &\text{die Werte aus ii) bilden eine Wahrscheinlichkeitsverteilung}.
\end{align*}
\end{ergebnis}

\begin{kontrolle}
$a$: für $n=2$: $1-4+4=1$; für $n=3$: $1-6+12-8=-1$. Passt. b): Für $n=2$,
$p=\tfrac12$ ist $\sum_{k=0}^{2}\binom2k\tfrac1{2^k}\tfrac1{2^{2-k}}
=\tfrac14(1+2+1)=1$.
\end{kontrolle}

\begin{merksatz}
Der binomische Lehrsatz liefert sowohl den Wert von $\sum\binom nk(-2)^k$
als auch den Nachweis, dass die einzelnen
``genau-$k$-Mutationen''-Wahrscheinlichkeiten sich zu $1$ summieren.
\end{merksatz}
